\section{Tìm vị trí phần tử lớn nhất [Cơ bản]}

\subsection{Phân tích \& Thiết kế (M0)}
\begin{itemize}
    \item \textbf{Input:}
    \begin{itemize}
        \item File \texttt{array.txt} chứa: số nguyên $n$ và mảng $n$ số nguyên.
    \end{itemize}

    \item \textbf{Xử lý:}
    \begin{itemize}
        \item \textbf{Case biên:}
        \begin{itemize}
            \item $n = 0$ $\rightarrow$ ghi \texttt{EMPTY}.
            \item $n = 1$ $\rightarrow$ \texttt{max\_val = a[0]}, \texttt{max\_idx = 0}.
        \end{itemize}

        \item \textbf{Thuật toán:}
        \begin{itemize}
            \item Khởi tạo \texttt{max\_val = a[0]}, \texttt{max\_idx = 0}.
            \item Duyệt mảng từ $i = 1$ đến $n-1$:
            \begin{itemize}
                \item Nếu \texttt{a[i] > max\_val} thì cập nhật.
            \end{itemize}
            \item Dùng dấu \texttt{>} để đảm bảo khi có nhiều phần tử lớn nhất bằng nhau thì lấy chỉ số nhỏ nhất.
        \end{itemize}

        \item \textbf{Module 1 (Đọc dữ liệu):} Hàm \texttt{readArray}. Bắt lỗi mở file (\texttt{return false}).
        \item \textbf{Module 2 (Xử lý \& Ghi kết quả):} Hàm \texttt{writeOutput} tìm max và ghi ra file.
    \end{itemize}

    \item \textbf{Output:}
    File \texttt{output.txt}:
\begin{lstlisting}
max_val=<i> / max_idx=<i>
\end{lstlisting}
\end{itemize}

\subsection{Mã nguồn C++11 (Tách module)}
\begin{lstlisting}[language=C++, breaklines=true]
    #include <iostream>
    #include <fstream>
    #include <vector>

    using namespace std;

    // MODULE 1: Doc mang tu file
    bool readArray(string path, vector<int>& a) {
        ifstream file(path);
        if (!file.is_open()) {
            cerr << "Khong mo duoc file " << path << "\n";
            return false;
        }

        int n;
        file >> n;

        int x;
        for (int i = 0; i < n; i++) {
            file >> x;
            a.push_back(x);
        }

        file.close();
        return true;
    }

    // MODULE 2: Tim max va ghi ket qua
    bool writeOutput(string path, const vector<int>& a) {
        ofstream file(path);
        if (!file.is_open()) {
            cerr << "Khong ghi duoc file " << path << "\n";
            return false;
        }

        if (a.empty()) {
            file << "EMPTY";
            file.close();
            return true;
        }

        int max_val = a[0];
        int max_idx = 0;

        for (int i = 1; i < (int)a.size(); i++) {
            if (a[i] > max_val) {
                max_val = a[i];
                max_idx = i+1;
            }
        }

        file << "max_val=" << max_val
            << " / max_idx=" << max_idx;

        file.close();
        return true;
    }

    // MAIN
    int main() {
        vector<int> a;

        if (!readArray("array.txt", a)) return 1;
        if (!writeOutput("output.txt", a)) return 1;

        return 0;
    }
\end{lstlisting}

\subsection{Kế hoạch kiểm thử (Test Cases)}

\textbf{Test Case 1 (Thông thường)}

\textbf{Input (array.txt):}
\begin{lstlisting}
5 1 5 3 9 2
\end{lstlisting}

\textbf{Expected Output (output.txt):}
\begin{lstlisting}
max_val=9 / max_idx=4
\end{lstlisting}

\vspace{0.5cm}

\textbf{Test Case 2 (Case biên -- nhiều phần tử lớn nhất bằng nhau)}

\textbf{Input (array.txt):}
\begin{lstlisting}
6 4 8 2 8 1 8
\end{lstlisting}

\textbf{Expected Output (output.txt):}
\begin{lstlisting}
max_val=8 / max_idx=2
\end{lstlisting}

\vspace{0.5cm}

\textbf{Test Case 3 (Case biên -- n = 0)}

\textbf{Input (array.txt):}
\begin{lstlisting}
0
\end{lstlisting}

\textbf{Expected Output (output.txt):}
\begin{lstlisting}
EMPTY
\end{lstlisting}

\vspace{0.5cm}

\textbf{Test Case 4 (Case biên -- n = 1)}

\textbf{Input (array.txt):}
\begin{lstlisting}
1 42
\end{lstlisting}

\textbf{Expected Output (output.txt):}
\begin{lstlisting}
max_val=42 / max_idx=1
\end{lstlisting}
